\section{引言与课程回顾}

本讲是 KAIST CS492D ``Diffusion Models and Their Applications'' 课程的第二讲，由 Minhyuk Sung 教授讲授。本讲从上一讲（Lecture 01）介绍的生成模型基础出发，逐步引出\textbf{去噪扩散概率模型}（Denoising Diffusion Probabilistic Models, DDPM）的核心思想。

在上一讲中，我们已经回顾了以下关键概念：
\begin{itemize}
    \item 给定数据集 $\{x^{(i)}\}$，如何将数据视为某未知概率分布 $p_{\text{data}}(x)$ 的样本；
    \item 生成模型的核心任务：学习一个从简单分布（通常是标准高斯 $\mathcal{N}(0, I)$）到数据分布的映射；
    \item GAN 和 VAE 两种生成模型范式的基本思路。
\end{itemize}

\begin{knowledgebox}{生成模型的统一框架}
所有生成模型的核心思想可以概括为：找到一个从简单的\textbf{先验分布}（prior distribution）$p(z) = \mathcal{N}(0, I)$ 到\textbf{数据分布}（data distribution）$p_{\text{data}}(x)$ 的映射。生成过程即为：先从 $p(z)$ 采样得到 $z$，再通过映射函数将 $z$ 变换为数据空间中的样本 $x$。
\end{knowledgebox}

\subsection{关键术语回顾}

在进入 DDPM 之前，需要牢记以下贝叶斯框架中的术语，它们贯穿整个课程：

\begin{table}[H]
\centering
\begin{tabular}{lll}
\toprule
\textbf{术语} & \textbf{符号} & \textbf{含义} \\
\midrule
先验分布（Prior） & $p(z)$ & 潜变量的分布，通常设为 $\mathcal{N}(0, I)$ \\
似然（Likelihood） & $p_\theta(x|z)$ & 给定潜变量生成数据的条件分布 \\
边际分布（Marginal） & $p_\theta(x)$ & 数据的边际概率 $\int p_\theta(x|z)p(z)dz$ \\
后验分布（Posterior） & $p_\theta(z|x)$ & 给定数据推断潜变量的分布 \\
\bottomrule
\end{tabular}
\caption{贝叶斯框架中的关键术语}
\end{table}

\subsection{本章小结}

本节回顾了生成模型的基本框架和关键术语，为后续引入 VAE 的 ELBO 推导和扩散模型的层级化扩展打下基础。


